\documentclass[twocolumn]{aastex631}
\begin{document}
\title{IllustrisTNG (TNG100-1, z=0) stellar mass function from 203,524 galaxies in a (111 Mpc)3 box; n(>1010.5M⊙)=1.49e-03 Mpc-3}
\author{NebulaMind Lab (autonomous pipeline)}
\affiliation{NebulaMind Lab, \url{https://lab.nebulamind.net}}
\begin{abstract}
IllustrisTNG (TNG100-1, z=0) stellar mass function from 203,524 galaxies in a (111 Mpc)3 box; n(>1010.5M⊙)=1.49e-03 Mpc-3. Generated autonomously from public data (tng) via the NebulaMind Lab runner: a bounded, reproducible, descriptive study.
\end{abstract}
\section{Introduction}
This short study was configured through the NebulaMind Lab and computed automatically. The requested analysis is a stellar-mass-function using TNG.
\section{Data and method}
Data are taken directly from tng. We form median relations in stellar-mass bins without further homogenisation of IMF or abundance calibration.
\begin{figure}\centering\includegraphics[width=\columnwidth]{result.png}\caption{IllustrisTNG (TNG100-1, z=0) stellar mass function from 203,524 galaxies in a (111 Mpc)3 box; n(>1010.5M⊙)=1.49e-03 Mpc-3}\end{figure}
\section{Result}
IllustrisTNG (TNG100-1, z=0) stellar mass function from 203,524 galaxies in a (111 Mpc)3 box; n(>1010.5M⊙)=1.49e-03 Mpc-3.
\section{Caveats}
This is an \emph{automated} first-pass descriptive result. It uses default selections and calibrations, does not homogenise IMF or abundance scales across sources, and applies no completeness or selection modelling. It is a starting point, not a validated measurement.
\end{document}
